% 2016 年全国卷Ⅱ理科第 8 题：秦九韶算法程序框图
\begin{tikzpicture}[>=stealth,line cap=round,line join=round,font=\small]
  % 开始与输入 x,n
  \draw[rounded corners=4pt] (-0.5,6.94) rectangle (0.5,7.36);
  \node at (0,7.15) {开始};
  \draw (-1.0,6.45)--(-0.8,6.0)--(0.8,6.0)--(1.0,6.45)--cycle;
  \node at (0,6.225) {输入 $x,n$};

  % 初始化、输入 a、循环计算
  \draw (-1.1,4.98) rectangle (1.1,5.48);
  \node at (0,5.23) {$k=0,\ s=0$};
  \draw (-0.82,4.47)--(-0.65,3.97)--(0.65,3.97)--(0.82,4.47)--cycle;
  \node at (0,4.22) {输入 $a$};
  \draw (-1.1,2.55) rectangle (1.1,3.35);
  \node at (0,2.95) {\shortstack{$s=s\cdot x+a$\\$k=k+1$}};

  % 判断、输出与结束
  \draw (0,2.25)--(1.3,1.8)--(0,1.35)--(-1.3,1.8)--cycle;
  \node at (0,1.8) {$k>n$};
  \draw (-0.82,0.8)--(-0.65,0.3)--(0.65,0.3)--(0.82,0.8)--cycle;
  \node at (0,0.55) {输出 $s$};
  \draw[rounded corners=4pt] (-0.5,-0.66) rectangle (0.5,-0.24);
  \node at (0,-0.45) {结束};

  % 主流程与“否”分支回到输入 a
  \draw[->] (0,6.94)--(0,6.45);
  \draw[->] (0,6.0)--(0,5.48);
  \draw (0,4.98)--(0,4.72);
  \draw[->] (0,4.72)--(0,4.47);
  \draw[->] (0,3.97)--(0,3.35);
  \draw[->] (0,2.55)--(0,2.2);
  \draw[->] (0,1.4)--(0,0.8);
  \draw[->] (0,0.3)--(0,-0.24);
  \draw (1.3,1.8)--(2.25,1.8)--(2.25,4.72)--(0,4.72);
  \node[right=2pt] at (1.3,2.0) {否};
  \node[right=3pt] at (0,1.08) {是};
\end{tikzpicture}
